\usection{Documentación del Lenguaje y Entorno}

Durante el proceso de desarrollo, con el propósito de garantizar la viabilidad, escalabilidad y correcta transferencia de la herramienta desarrollada, se elaboró una documentación conformada por dos recursos principales: el manual de sistema y el manual de usuario. El manual de sistema está orientado estrictamente a desarrolladores y personal técnico, detallando la arquitectura del motor de ejecución concurrente, la semántica del flujo de datos y la estructuración de los protocolos de comunicación distribuida (WebSockets y protocolo para las tarjetas micro:bit). Este documento resulta indispensable para facilitar tareas de mantenimiento, resolución de problemas técnicos y el futuro desarrollo de nuevas extensiones o nodos para el lenguaje. El documento puede hallarse en el apéndice~\ref{apx:system-manual}.

Por su parte, el manual de usuario fue concebido como una guía didáctica dirigida tanto a los estudiantes como a los educadores. Este documento es de particular interés para el proyecto, ya que trasciende la simple descripción de la interfaz y el catálogo de nodos al incorporar una serie de retos simples y prácticos que demuestran, paso a paso, cómo utilizar el sistema para construir algoritmos funcionales.

La estructuración de este manual a través de ejemplos guiados es de fundamental importancia para dotar a los docentes de un recurso claro que les permita conocer a fondo la herramienta y aplicarla con seguridad en sus aulas de clase, dando así respuesta directa a la recomendación formativa emitida por la docente experta durante la etapa de validación pedagógica. El manual se integró en la página web, facilitando su acceso, y puede verse en el apéndice \ref{apx:user-manual}.
